%root = ../thesis-main.tex
\chapter{Conclusions}\label{ch:conclusions}
\section{Summary of Contributions}

This thesis introduced a fully reactive approach to discrete-event simulation.
It begins by establishing the theoretical foundations of system modeling and
simulation and conducting an in-depth analysis of reactive programming
paradigms and frameworks. The work culminates in the implementation of a
working prototype that reframes dependency management, where instead of
constructing or repairing explicit dependency graphs, the problem is redefined
as maintaining consistent data flows. Accordingly, the central contribution is
a reactive architectural perspective in which propensities and rates are
treated as derived values. These values update automatically as soon as their
underlying state changes, thereby reducing complexity in updates management
while preserving the eager consistency expected by a discrete-event scheduler.

On this basis, a lightweight reactive framework tailored to Alchemist was
designed and implemented around a minimal \kotlin{Observable<T>} contract. This
framework addresses lifecycle concerns through explicit subscriptions and a
rigorous disposability pattern, ensuring that resources for out-of-scope
elements are released by design rather than by convention. This foundation
enables a compositional style for modelling simulation dynamics, since
monadic operators such as \kotlin{map}, \kotlin{flatMap} and
\kotlin{switchMap}, together with higher level combinators, allow to
express complex dependencies as chains of small transformations, thus improving
modularity and reducing boilerplate.

Furthermore, the work introduced reactive collections which support both
structural changes and fine grained element updates, and consequently allow
aggregate signals to remain consistent even as neighbourhoods evolve and local
state changes propagate through dynamically assembled computations. Crucially,
the proposed framework optimises derived value representation to both minimise
the memory overhead often associated with naive reactive implementations and
introduce the framework to laziness opportunities. The successful replacement
of the Alchemist engine with its reactive counterpart demonstrates
\begin{enumerate*}[label=(\roman*)]
	\item the applicability of reactive programming principles to
	      discrete-event simulation,
	\item the expressive power gained from compositional reactive primitives
	      and
	\item the theoretical performance gains achievable in scenarios where
	      component updates are sparse and loosely coupled.
\end{enumerate*}

\section{Limitations and Trade-offs}
\label{sec:limitations}

While the proposed reactive architecture successfully addresses the maintenance
complexity of the dependency graph and can theoretically improve performance in
sparse scenarios, it is not without significant trade-offs. The shift from an
explicit, centralised graph to a decentralised, event-driven model introduces
specific limitations regarding memory consumption and pressure, performance in
dense topologies and code comprehension and debugging complexity. Finally, even
if the fully-reactive event driven simulator has been proven to be correct, it
still have no means of achieving glitch freedom.

\subsection{Memory Footprint and Object Overhead}
In the classic Alchemist engine, a dependency is represented by a
lightweight edge in a graph. In the reactive framework, a dependency is
materialised as a chain of objects:
\begin{itemize}
	\item Each \kotlin{Observable} maintains a collection of subscribers (e.g.
	      a \kotlin{LinkedHashMap} and \kotlin{ArrayList}).
	\item Each transformation (e.g. \kotlin{map}, \kotlin{filter}) creates a
	      new intermediate \kotlin{DeivedObservable} instance.
	\item Crucially, every subscription involves the creation of closure
	      objects (lambdas) to capture the context, effectively doubling the
	      object count for every link in the dependency chain.
\end{itemize}
Even if results in \Cref{ssec:benchmark} shows similar usage compared to the
original implementation, in scenarios with millions of molecules and reactions,
the sheer volume of \kotlin{Observable} wrappers and listener objects can
pressure the \ac{JVM} garbage collector, leading to a possibly increasing
performance degradation due to the high volume of memory operations.

\subsection{Performance Penalty in Dense Graphs}
The reactive model thrives on \emph{sparse} scenarios, where a change in one
component affects only a few others. However, in \emph{dense} or \emph{highly
	interdependent} scenarios, the overhead of the reactive propagation mechanism
outweighs its benefits.

In the standard engine, updating a dense dependency graph involves iterating
over a known list of edges, a linear operation with high cache locality. In the
reactive engine, a single state change triggers a cascade of lookups, lambda
invocations and context switches through following the \kotlin{onChange}
chains.

This function-call overhead is negligible for shallow chains but accumulates
rapidly in deep dependency trees (e.g. the multi-step protein interactions in
the Repressilator). Consequently, for simulations where the global state
changes significantly at every step (i.e. dense dependency graphs), the classic
explicit-graph approach remains more efficient.

\subsection{Inversion of Control and Debuggability issues}
A well-known challenge in reactive programming is the set of problems derived
by the inversion of control. In the original engine, the flow of execution is
explicit: the engine selects a reaction, checks its dependencies, and executes
it. This linear flow is familiar to most programmers and easy to trace with a
standard debugger.

In contrast, the reactive architecture inverts parts of this flow. As argued in
\cite{maier2010DOP}, this inversion forces developers to introduce side effects
to express local state mutation upon the receipt of an event. This often
violates encapsulation, leading to designs that are generally harder to reason
about compared to their reactive counterpart. Moreover, Inverting the control
flow typically causes a \emph{semantic segmentation}. In the reactive engine
the logic for a single reaction execution is no longer contained withing a
single, contiguous block of code, but instead it is fragmented across possibly
several asynchronous-like callbacks and intermediate operators.

This fragmentation has severe implications when it comes to debugging. A
breakpoint placed in a reaction's \kotlin{update} method does not reveal
\emph{why} it was triggered. To complicate matters even more, the call stack is
dominated by generic infrastructure code (e.g. \kotlin{DerivedObservable}'s
methods) rather than domain-specific logic. The obfuscation of the \emph{causal
	chain} makes interaction bugs harder to spot, often requiring specialised
visualisation tools or verbose logging to trace propagation of events through
the system.

\subsection{Transient Inconsistencies and Computational Waste}
A strict requirement for many reactive frameworks is \emph{glitch freedom},
namely the guarantee that no observer ever sees a stale or inconsistent state
during the propagation of an update (discussed in \Cref{ssec:glitches}).

Alchemist's reactive architecture does not strictly enforce glitch freedom in
the topological sense. When a single logical event triggers multiple updates
(e.g. a node moving impacts the distance to multiple neighbors), a dependent
reaction may receive multiple notifications in rapid succession. The first
notification might reflect an intermediate, inconsistent state, potentially
triggering a re-evaluation of the reaction's propensity based on this transient
data.

However, unlike systems where glitches can cause multiple modification of
global, shared variables, in the presented architecture these transient
inconsistencies do not corrupt the simulation correctness. The Alchemist engine
in fact, acts as a consistency gatekeeper through its update logic, where
reactions are effectively rescheduled if their newly computed scheduled time
($\tau$) is different from the previous value. This filters out transient
fluctuations caused by glitches, ensuring that only the final, consistent state
affects the simulation trajectory.

While correctness is preserved, the \emph{cost} of these glitches is
non-negligible. Every intermediate notification triggers a full re-computation
of the reaction's conditions and rate equation. In dense dependency graphs
where a single root change fans out to many intermediates before converging,
this results in \emph{redundant computation}. The reaction might be
re-evaluated $N$ times for a single logical state change, wasting CPU cycles on
$N-1$ invalid states. This computational waste is the primary trade-off
accepted to avoid the complexity of a topological sort propagation engine, and
will be subject of future works..

\subsection{Initialisation Fragility}
Finally, the presence of the \emph{set-up phase} and the \emph{execution phase}
remains a fragile point in this reactive architecture. As discussed in
\Cref{par:simulation_phases}, creating observables requires a stable
environment. Wiring reactions during the initialisation phase, i.e. before all
nodes are fully deployed, can lead to race conditions or ``dead'' observables
if not handled with strict discipline. While the framework provides mechanisms
to mitigate this, it shifts the burden of correctness onto the model developer,
who must now be acutely aware of the lifecycle state of the components they are
observing.

\section{Future Works}
The research presented in this thesis characterizes the reactive architecture
as a high-potential prototype. While the current implementation demonstrates
feasibility and theoretical performance gains in specific scenarios, the
analysis of its limitations outlines a clear roadmap for future development.
Future work will focus on optimising the core engine and enhancing the
developer experience of the reactive framework.

The primary objective is to mitigate the overhead observed in dense dependency
graphs. Key directions include:
\begin{itemize}
	\item \emph{Incarnation}-specific optimisations should be carried out, with
	      a deeper focus on the SAPERE project and its finer-grained dependencies
	      and complex template matching logic.
	\item Transitioning the Alchemist standard library (Reactions, Conditions)
	      from the current hybrid adapter pattern to a \emph{native reactive
		      implementation}. Removing the polling bridges would eliminate the
	      synchronization overhead and fully exploit the push-based paradigm;
	\item A comprehensive study on the \emph{resource usage} of the reactive
	      primitives. Future iterations could employ \emph{object pooling} or
	      specialized primitive streams (avoiding boxing) to reduce the pressure
	      on the \ac{GC}, thereby stabilizing performance in long-running
	      simulations;
	\item Introducing a \emph{transaction manager} to support atomic state
	      updates. This is crucial not only for glitch freedom but for
	      performance: by coalescing multiple updates into a single notification,
	      the engine can eliminate the redundant computations caused by transient
	      intermediate states. Moreover, the capability of performing batch
	      modifications atomically is added to the domain, where updates to
	      dependencies are propagated only at the end of the transaction.
\end{itemize}

Similarly, the \emph{framework} itself could undergo a series of revision to
enhance both its expressive power and its safety by design. More in particular,
drawing inspiration from \textsc{Scala.React}~\cite{maier2010DOP} and project
Reactor, the framework could leverage Kotlin's advanced features (e.g. context
receivers, internal DSLs, property delegates) to mitigate the inversion of
control problem, enforcing encapsulation and separation of concern by design by
using reactors and an imperative programming style in a reactive environment.
Moreover, the developer experience could be furthermore enhanced by providing
an out-of-the box debug toolbox, for example by visualising the dependency
graph at runtime to let developers investigate and trace the causal propagation
of events, solving the ``obscured causal chain'' issue.

